\section*{Chapter 16}
\paragraph{16.A1.} Low-rank signal contributes a few large singular values;
spread-out noise contributes a longer tail.
\paragraph{16.A2.} Count the entries in $L$ and $R$ and compare with $mn$.
\paragraph{16.B1.} Look for a large drop in the singular values.
\paragraph{16.B2.} Compare $3(100+80)$ with $100\cdot 80$.
\paragraph{16.C1.} Build a rank-$2$ signal as \texttt{L @ R} with two
columns/rows.
\paragraph{16.C2.} Use \texttt{np.cumsum(s**2) / np.sum(s**2)}.
\paragraph{16.D1.} Think of discarded singular directions as removed content
and retained tail directions as possible noise.
